% Methods

We compare four VQC architecture families, all built on the
\texttt{StronglyEntanglingLayers} ansatz from PennyLane as the base
quantum circuit block.

\subsection{Multi-Layer Fully-Connected VQC (FC-VQC)}

The FC-VQC cascades multiple VQC blocks with decreasing qubit counts,
connected via classical fully-connected (dense) layers.  For example,
the notation $16\mathrm{t}4\mathrm{t}1$ denotes a three-block architecture
mapping 16~qubits $\to$ 4~qubits $\to$ 1~qubit, with each block consisting
of \texttt{StronglyEntanglingLayers} at a fixed depth $d$ (typically~3).

Two inter-block connectivity modes are supported:
\begin{itemize}
  \item \textbf{Type~4 (fully-connected):} Each qubit measurement from
    block $\ell$ is connected to all angle-encoding inputs of block
    $\ell{+}1$ via a learned linear projection.
  \item \textbf{Type~3 (circular-shift):} A fixed permutation maps
    block $\ell$ outputs to block $\ell{+}1$ inputs, reducing parameters.
\end{itemize}

\subsection{Residual VQC (ResNet-VQC)}

ResNet-VQC augments the FC-VQC topology with classical residual connections.
When the qubit count is preserved between adjacent blocks ($n_\ell = n_{\ell+1}$),
a skip connection adds the input measurements to the output:
\begin{equation}
  \mathbf{z}_{\ell+1} = \mathrm{VQC}_{\ell+1}(\mathbf{z}_\ell) + \mathbf{z}_\ell.
\end{equation}
When dimensions differ, a learned linear projection is applied to the skip
path.  This mirrors classical ResNets~\cite{he2016deep} and is motivated by
the same goal: easing gradient flow through deep stacks of nonlinear blocks.

\subsection{Quantum Transformer VQC (QT --- Route A)}

QT implements a hybrid quantum--classical transformer block.  The input
features are first encoded into a VQC, whose measurement outputs serve as
token embeddings.  A \emph{classical} multi-head self-attention layer
operates on these embeddings, followed by a VQC-based feed-forward network:
\begin{equation}
  \mathbf{h} = \mathrm{Attention}(\mathrm{VQC}_{\mathrm{enc}}(\mathbf{x}))
  + \mathrm{VQC}_{\mathrm{enc}}(\mathbf{x}),
\end{equation}
\begin{equation}
  \mathbf{y} = \mathrm{VQC}_{\mathrm{ffn}}(\mathrm{LayerNorm}(\mathbf{h}))
  + \mathbf{h}.
\end{equation}
This route uses classical attention but quantum encoding and feed-forward
processing.

\subsection{Fully Quantum Transformer VQC (FQT --- Route B)}

FQT replaces the classical attention with a parameterized quantum circuit
that computes query--key interactions via entangling gates.  Both the
attention mechanism and the feed-forward network are implemented as VQCs,
making the entire transformer block quantum (except for measurement and
re-encoding).

We also study a \textbf{+LayerNorm} variant that applies classical
LayerNorm to the FQT outputs before the next layer, which our ablation
shows is important for classification performance.

\subsection{Noise Model}
\label{sec:noise-model}

To assess near-term hardware viability, we inject single-qubit depolarizing
noise after every layer of \texttt{StronglyEntanglingLayers}:
\begin{equation}
  \mathcal{E}(\rho) = (1 - p_d)\,\rho + \frac{p_d}{3}(X\rho X + Y\rho Y + Z\rho Z),
\end{equation}
where $p_d \in \{0.001, 0.005, 0.01, 0.05\}$ controls the noise strength.
% TODO: Expand after noise experiments are complete.
